\section{Introduction}\label{sec:intro}

Global structure-from-motion (SfM) estimates all camera rotations at once from pairwise relative rotations, then all positions, and finally refines everything with bundle adjustment. It is the backbone of GLOMAP~\cite{pan2024glomap}, now shipped inside COLMAP~\cite{schoenberger2016sfm}, and of recent hybrids that feed local reconstructions from networks such as VGGT~\cite{wang2025vggt} and MASt3R~\cite{leroy2024mast3r,duisterhof2025mast3rsfm} into global averaging~\cite{pan2026gluemap}. Rotation averaging is the first and most fragile stage: a wrong rotation propagates into translation averaging and bundle adjustment, and symmetric or repetitive scenes produce relative rotations that are wrong yet look correct~\cite{zach2010disambiguating,heinly2014correcting,cai2023doppelgangers}.

Two developments make this fragility more pressing. First, front ends now report uncertainty, and averaging methods can use it~\cite{zhang2023revisiting,olsson2025certifiably}. These methods assume the reported uncertainty is right. For learned front ends it is not: a 2026 audit of seven feed-forward reconstruction models found that confidences rank errors moderately well (median Spearman correlation $0.46$) but are overconfident by a median factor of $2.4$, with nominal 90\% intervals covering only 17--41\% of errors~\cite{deng2026audit}. Second, the level of the error changes with the data: fog, rain, low light or a new sensor raise the error at a fixed reported confidence. A pipeline that trusts the confidence, or a threshold tuned on clean data, then fails silently.

Robust rotation averaging, typically solved by iteratively reweighted least squares (IRLS), handles outliers well when the noise level matches the kernel it was tuned for~\cite{chatterjee2018robust,hartley2011l1,sidhartha2021weights}. It gives no statement about which cameras can be trusted, and it does not notice when the noise level has changed. Conformal prediction~\cite{vovk2005algorithmic,angelopoulos2023gentle} and conformal risk control (CRC)~\cite{angelopoulos2024crc} give distribution-free guarantees, and they have been applied to object pose~\cite{yang2023object}, visual odometry~\cite{stutts2023conformal}, single camera poses~\cite{pathak2026adaptive}, and failure gating for feed-forward reconstruction~\cite{ahmedt2026geocond}. GeoCond~\cite{ahmedt2026geocond} already gates view-graph edges with calibrated risk thresholds derived from learned per-pair features. To our knowledge, no prior work calibrates a cycle-closure test on the view graph by conformal risk control over scenes, gives coverage guarantees for the cameras of a jointly estimated global reconstruction, or detects from the scene itself that the front end's uncertainty has drifted.

\paragraph*{Contributions} We present CARA (Calibrated, conformally gated, drift-Aware Rotation Averaging), a rotation-averaging method for global SfM that returns rotations together with an edge acceptance set, a registration set, and a per-camera error bound. Our contributions are:
\begin{enumerate}
\item A \emph{cycle nonconformity score} that normalizes loop-closure error by recalibrated front-end uncertainty, and \emph{conformal edge and residual gates} whose thresholds are set by conformal risk control over scenes, each with a finite-sample bound on the expected fraction of inlier edges it rejects (\cref{thm:edge,thm:res}).
\item A \emph{label-free drift factor} $\hat\rho$ estimated from cycle closure on the test scene, with a population bracket under outlier contamination (\cref{prop:bracket}), and a small-noise equivariance result showing that rescaling uncertainties and the robust kernel by $\hat\rho$ restores the calibrated behavior under a pure scale shift (\cref{lem:equiv}).
\item \emph{Per-camera error certificates} for a jointly estimated reconstruction, with a finite-sample bound on the expected fraction of registered cameras whose error exceeds its bound (\cref{thm:cam}), together with \emph{consensus bridging}, which lets cameras abstain instead of being attached through unverified edges.
\item \emph{Extensions for camera-level failures}, which cycle tests cannot detect: calibration on scenes with modeled failures, difficulty-conditional bounds with a group-wise guarantee (\cref{cor:group}), and an edge uncertainty learned from several front-end signals, tested in a pre-registered confirmatory run together with their effect on camera positions after translation averaging.
\item An evaluation that reports failures as well as gains: ten-seed synthetic studies with seed-level significance tests, ablations of the main components (\supp{2}), a structured-doppelganger study where cycle tests fail by construction, real data from the 1DSfM Notre Dame collection~\cite{wilson2014robust} with a strictly disjoint calibration/test split and controlled shifts, and real images from ETH3D~\cite{schops2017eth3d} with the VGGT front end under rendered fog, low light and rain, calibrated and tested on different places.
\end{enumerate}

\paragraph*{Scope} CARA addresses the rotation stage. Translation averaging, bundle adjustment and a full comparison with GLOMAP on large image collections are outside the scope of this paper; \cref{sec:discussion} states what they would require. The code, data splits and all result logs are publicly available (see the data availability statement).
